% Material parked from v1 (../02_fixed-h-mcfg/main.tex) when building 02-v2, 2026-10-30.
% Not part of v2. Kept for possible use in a later paper or the thesis.
% (2) v1 Proposition prop:unary-cyclic-linear (one-letter cyclic typing, tau <= 2m theta-bar) and its proof.

% (1) restored to main.tex as an appendix on 2026-10-30 (after the ChatGPT comparison review).

%% ---- (2) ----
\begin{proposition}\label{prop:unary-cyclic-linear}
Let $\Sigma=\{\mathtt a\}$, $m\ge2$, and let $s$ send $\mathtt a$ to a generator of
$\mathbb Z_m$.  For every $L\in\Cmcf{f}{s}$ and every reduced good
fan-out-$\le f$, rank-$\le r$ presentation $G$ of $L\cap\Sigma^+$,
$\tau(G,s)\le2m\bar\theta_G$.  Hence $\mathcal B^{\le r}_{f,s}$ and
$\mathcal B_{f,s}$ have a common characteristic sample $K$ with
$\|K\|_+\le(m+1)^{fr}(3+2mr)\|G\|\bar\theta_G+\mathbf 1_{\{\eps\in L\}}$.
\end{proposition}

\begin{proof}
Let $A:=\{n\ge1:\mathtt a^n\in L\}$, assumed nonempty, $\theta:=\bar\theta_G$ and
$\mu:=\min A\le\theta$.  (a) If $n<n'$ lie in $A$ and $m\mid\delta:=n'-n$, then
$\mathtt a$ and $\mathtt a^{1+\delta}$ have the same type and share the context
$(\mathtt a^{n-1},\eps)$, so testing their common distribution on
$(\mathtt a^{p-1},\eps)$ gives $p\in A\iff p+\delta\in A$ for all $p\ge1$.

If $A$ is finite, its elements are pairwise incongruent modulo $m$ by (a), so
$|A|\le m$, and Lemma~\ref{lem:app-lf-descent} reaches $\mu$ from $\max A$ in
at most $|A|-1$ steps of length at most $\theta-\mu$; hence $\max A\le m\theta$.
Every typed yield embeds into a word of $L$, so $\tau\le t_{\mathrm{typ}}\le m\theta$.

Let $A$ be infinite and let $\Delta$ be the greatest common divisor of the set $D$ of
differences $n'-n>0$ of elements of $A$ congruent modulo $m$; then $m\mid\Delta$.
Fix $\delta_0\in D$ and let $\chi(z)=1$ iff $z+N\delta_0\in A$ for some, equivalently
by (a) every, $N$ with $z+N\delta_0\ge1$; so $\chi$ extends the indicator of $A$ to
$\mathbb Z$.  For $\delta\in D$ and $z\in\mathbb Z$, (a) applied to $p=z+N\delta_0$ with $N$
large gives $\chi(z+\delta)=\chi(z)$.  The periods of $\chi$ form a subgroup of
$\mathbb Z$, so $\Delta$ is a period.  Two elements of $A$ congruent modulo $m$ differ by a multiple of $\Delta$, so
$A$ meets at most $m$ residue classes modulo $\Delta$, and some consecutive elements
$u<v$ of $A$ with $u\ge\mu$ satisfy $v-u\ge\Delta/m$.  Lemma~\ref{lem:app-lf-descent}
gives an element of $A$ in $[v-\theta,v)$, hence at most $u$, so $\Delta\le m\theta$.
Now let $X$ be a typed nonterminal of arity $d$ and $\tuple x\in L_X$ of total
length $T$.  Choose $\rho_h\in\{1,\ldots,m\}$ with $\rho_h\equiv|x_h|\pmod m$ and
put $b:=\sum_h\rho_h\le md\le m\theta$, using $d\le\|\omega\|_1\le\theta$ for a short
tuple $\omega$ of the underlying nonterminal.  Since $T\equiv b\pmod m$ and
$m\mid\Delta$, some $T'\equiv T\pmod\Delta$ satisfies $b\le T'<b+\Delta$ and $m\mid T'-b$.  The
tuple $\tuple y:=(\mathtt a^{\rho_1+T'-b},\mathtt a^{\rho_2},\ldots,\mathtt a^{\rho_d})$
has the type of $\tuple x$, and $|F[\tuple y]|\equiv|F[\tuple x]|\pmod\Delta$ for every
$d$-context $F$, so $\tuple x$ and $\tuple y$ have the same distribution.  Hence
$\tau_X\le T'<b+\Delta\le2m\theta$.

Corollary~\ref{cor:distribution-representatives} gives a common characteristic
sample of length at most $N_{\mathrm{typ}}(1+2\theta+r\tau)+\mathbf 1_{\{\eps\in L\}}$,
and $N_{\mathrm{typ}}\le|V|(m+1)^f+|P|(m+1)^{fr}\le(m+1)^{fr}\|G\|$ since
$|\mathsf S^\sharp|=m+1$.
\end{proof}

